\hwhat{In 14 non-holdout goal periods, a 5-min hazard model (agent, project, day effects) tests whether a seen chat link to project X raises switches to X. A seen link goes with a $2.3\times$ switch hazard (pooled; $p<0.05$ in 13/14). It beats links shifted by $\pm$30--120 min in only 3/11 herding weeks: the primary test fails. Future links predict switches better than past ones in 10/11. Naive agents already switch at $3.3\times$ baseline in the hour before their first exposure. Links lead visits only in own-artifact weeks, where nobody herds. Round 1b uses context-ledger visibility and work commits. The pre-NE09 weeks were mismeasured, so 5/11 herding weeks now beat the shift (primary test mixed; pooled $2.6\times$). Future links still win (9/11). Recipients switch at $6.5\times$ baseline in the $\sim$17\,s before they can read the link. NE09 was not a visibility change. Work switches show the same pattern. The holdout is not run.}

\hmodel{Multi-label kinetic Potts (10) with a contagion field (03). An off-project agent $i$ switches into $X$ at rate $\mu_{iX}$. $\alpha_i,\pi_X,\delta_d$: agent, project and day fields. $E_{iX}$: a link to $X$ by another agent became visible to $i$ in the last 60 min. $\kappa$: link coupling. $n_{-i,X}$: others on $X$ (Potts coupling $J$). $\mathbf h$: own history, plans, pulses. $\lambda$: extra switches per exposure. $R_{\rm link}$: share of switches attributable to links. $\pi$: links per switch. $k_s$: susceptible recipients per link.}

\hmath{\vspace{-5pt}\begin{gather*}
\log\mu_{iX}(t)=\alpha_i+\pi_X+\delta_d+\kappa\,E_{iX}(t)\\ +\,J\log\!\big(1+n_{-i,X}(t)\big)+\boldsymbol\eta\!\cdot\!\mathbf h_{iX}(t)\\
\lambda=\frac{\sum(\hat\mu-\hat\mu|_{E=0})}{N_{\rm exp}},\qquad R_{\rm link}=\pi\,k_s\,\lambda=\frac{\sum(\hat\mu-\hat\mu|_{E=0})}{\sum y}
\end{gather*}\vspace{-7pt}}

\hobs{../figures/summary_obs.pdf}{Round 1. (a) Switch-hazard coefficient for a link seen in the last hour (blue) and for a link that arrives in the next hour (orange), fitted jointly. In every herding week except \#24, the future link predicts as well or better. Only the own-artifact weeks (\#39, \#40, \#42) show past $>$ future. (b) Simulated peak number of agents on one project with half (green) or all (orange) links removed: median $\times0.79$ with no links, an upper bound.}

\htelos{Throttling link-sharing does not stop pile-ons. Removing every link trims simulated peak pile-ons by at most about 30\% (median $\times0.71$ in round 1b; an upper bound). In herding weeks, links, chat and work on a project rise together: the link marks a burst already under way. To prevent duplicated work, assign ownership or act on the burst triggers, not on links. Use the lead/lag test as a diagnostic. Lag $\gg$ lead means directed recruitment (agents advertising their own artifacts, $\lambda\approx0.01$--$0.02$ per exposure); lead $\ge$ lag means markers.}

\hbottom{Links co-move with project switches (HR $\approx2.6$) but do not lead them in herding weeks; they mark bursts, and removing them trims peaks by $\le$30\%.}

\hresults{\scriptsize\setlength{\tabcolsep}{2pt}\begin{tabular}{@{}p{0.30\columnwidth}p{0.50\columnwidth}p{0.17\columnwidth}@{}}
\textbf{Prediction} & \textbf{Observed (11 herding weeks)} & \textbf{Verdict}\\\hline
\raggedright P1 $\kappa>0$, $p<.05$, beats link shift & \raggedright 3/11 (\#24, \#30, \#51); pooled $e^\kappa=2.3$ & \raggedright failed\tabularnewline
\raggedright P2a past $>$ future links & \raggedright future $>$ past in 10/11; pooled $-1.17\pm0.30$ & \raggedright reversed\tabularnewline
\raggedright P2b survives momentum & \raggedright 3/3 & \raggedright yes\tabularnewline
\raggedright P3 naive agents; no pre-trend & \raggedright naive $\kappa=1.31$; pre-trend $3.3\times$ & \raggedright half\tabularnewline
\raggedright P4 simple contagion & \raggedright 1 link: HR 2.3; $\ge$2 senders add $-0.24$ & \raggedright yes\tabularnewline
\raggedright P5/P6 $\lambda$, $R_{\rm link}<0.5$ & \raggedright $\lambda$ median 0.036; $R_{\rm link}$ 0.06--0.44 & \raggedright yes (bounds)\tabularnewline
\raggedright P7 no links: peaks $-$10--30\% & \raggedright median $\times0.79$; calibration 9/14 & \raggedright mixed\tabularnewline
\raggedright P8 NE09 slows response & \raggedright pre-NE09 as fast; 1b: premise void (17 vs 18\,s delay) & \raggedright failed\tabularnewline
\raggedright 1b P1 on ledger visibility & \raggedright 5/11; pooled $e^\kappa=2.6$; future $>$ past $-1.04\pm0.24$ (9/11) & \raggedright mixed\tabularnewline
\raggedright 1b blind window (before reading) & \raggedright $6.5\times$ baseline vs $3.5\times$ after reading & \raggedright co-burst\tabularnewline
\raggedright 1b \#31 wave in work commits & \raggedright $\kappa$ +0.57, shift $z$ +0.2, future $>$ past & \raggedright failed\tabularnewline
\raggedright P9 own-artifact $\lambda$ lower & \raggedright 0.012--0.019 vs 0.036; past $>$ future & \raggedright yes\tabularnewline
\raggedright P10 occupancy $J>0$ & \raggedright 3/11 under both FE & \raggedright failed\tabularnewline
\end{tabular}}

\hobsb{../figures/synthetic_compact.pdf}{Synthetic swarms at village size, 25 runs each. Blue: the primary test calls contagion only when links truly act (92--100\%; 52\% with drive). Orange: future links beating past links is typical of drive, occupancy and announcement worlds (88--100\%) and rare under pure contagion (8--16\%). Real herding weeks show it in 10/11 (round 1).}

\hcaveats{Touches measure attention (strict artifact mentions); round 1b adds work commits in git-dense weeks only. Blind windows on the ledger are short: 114 blind switches in all. $\lambda$ and $R_{\rm link}$ are attributable associations, so they bound any causal link effect from above. The simulator has $\pm0.1$ error in synthetic tests, overstates the effect when drive and contagion mix, and under-predicts peaks in long periods (\#38, \#51). The other-room placebo is uninformative, because rooms work on different projects. On shared window labels no period beats the shift null; this check, the short-lag check, label robustness and the H27 link are post hoc.}

\hnext{Run the frozen test on \#22, \#28, \#45. Model what starts attention bursts (human messages, kickoffs, announcements), jointly with the H27 onsets. Split links by sender type and by announcement vs request (Phase 2 text).}
